\documentclass[10pt]{article}
\usepackage[a4paper, margin=0.75in]{geometry}
\usepackage{booktabs}
\usepackage{array}
\usepackage{longtable}
\usepackage{multirow}
\usepackage{colortbl}
\usepackage{xcolor}
\usepackage{caption}
\usepackage{graphicx}
\usepackage{amsmath}
\usepackage{siunitx}

% Times New Roman font settings
\usepackage{mathptmx} % Times font for text and math
\usepackage[T1]{fontenc}
\usepackage[utf8]{inputenc}

% Define colors for IEEE standard
\definecolor{headercolor}{HTML}{E8F4F8}
\definecolor{altrowcolor}{HTML}{F5F5F5}
\definecolor{accentcolor}{HTML}{2E86AB}

\begin{document}

\title{Age Distribution Analysis of Patient Population\\{\large Comprehensive Demographic Study}}
\author{}
\date{}
\maketitle

\section{Age Distribution Overview}

This analysis presents the age distribution patterns across 331 patients from the medical dataset, providing insights into demographic characteristics and age-related healthcare utilization patterns.

\section{Age Distribution Data Table}

\begin{table}[h!]
\centering
\caption{Patient Age Distribution by Decade Groups}
\begin{tabular}{c|c|c|c|c}
\toprule
\rowcolor{headercolor}
\textbf{Age Group} & \textbf{Number of} & \textbf{Percentage} & \textbf{Cumulative} & \textbf{Cumulative} \\
\rowcolor{headercolor}
\textbf{(years)} & \textbf{Patients} & \textbf{(\%)} & \textbf{Count} & \textbf{Percentage (\%)} \\
\midrule
0--10 & 0 & 0.0 & 0 & 0.0 \\
\rowcolor{altrowcolor}
11--20 & 11 & 3.3 & 11 & 3.3 \\
21--30 & 40 & 12.1 & 51 & 15.4 \\
\rowcolor{altrowcolor}
31--40 & 57 & 17.2 & 108 & 32.6 \\
41--50 & 83 & 25.1 & 191 & 57.7 \\
\rowcolor{altrowcolor}
51--60 & 78 & 23.6 & 269 & 81.3 \\
61--70 & 41 & 12.4 & 310 & 93.7 \\
\rowcolor{altrowcolor}
71--80 & 17 & 5.1 & 327 & 98.8 \\
81--90 & 3 & 0.9 & 330 & 99.7 \\
\rowcolor{altrowcolor}
91--100 & 1 & 0.3 & 331 & 100.0 \\
\midrule
\rowcolor{headercolor}
\textbf{Total} & \textbf{331} & \textbf{100.0} & \textbf{331} & \textbf{100.0} \\
\bottomrule
\end{tabular}
\end{table}

\section{Statistical Summary}

\subsection{Descriptive Statistics}

\begin{table}[h!]
\centering
\caption{Age Distribution Statistical Measures}
\begin{tabular}{l|c}
\toprule
\rowcolor{headercolor}
\textbf{Statistical Measure} & \textbf{Value} \\
\midrule
Total Patients & 331 \\
\rowcolor{altrowcolor}
Modal Age Group & 41--50 years (83 patients) \\
Peak Age Range & 41--60 years (161 patients, 48.7\%) \\
\rowcolor{altrowcolor}
Young Adults (21--40) & 97 patients (29.3\%) \\
Middle-aged (41--60) & 161 patients (48.7\%) \\
\rowcolor{altrowcolor}
Elderly (61+) & 62 patients (18.7\%) \\
Pediatric (0--20) & 11 patients (3.3\%) \\
\bottomrule
\end{tabular}
\end{table}

\subsection{Age Group Categories}

\begin{table}[h!]
\centering
\caption{Consolidated Age Categories for Clinical Analysis}
\begin{tabular}{l|c|c|l}
\toprule
\rowcolor{headercolor}
\textbf{Age Category} & \textbf{Age Range} & \textbf{Count} & \textbf{Percentage} \\
\midrule
Pediatric \& Adolescent & 0--20 years & 11 & 3.3\% \\
\rowcolor{altrowcolor}
Young Adult & 21--40 years & 97 & 29.3\% \\
Middle-aged & 41--60 years & 161 & 48.7\% \\
\rowcolor{altrowcolor}
Elderly & 61--80 years & 58 & 17.5\% \\
Very Elderly & 81+ years & 4 & 1.2\% \\
\bottomrule
\end{tabular}
\end{table}

\section{Key Demographic Insights}

\subsection{Age Distribution Patterns}

\begin{itemize}
\item \textbf{Peak Demographics:} The largest patient group falls within the 41--50 age range (83 patients, 25.1\%)
\item \textbf{Working Age Dominance:} Combined 31--60 age groups represent 218 patients (65.9\% of total population)
\item \textbf{Middle-age Concentration:} 41--60 age groups account for 161 patients (48.7\% of total)
\item \textbf{Limited Pediatric Representation:} Only 11 patients (3.3\%) under age 20
\item \textbf{Elderly Care Needs:} 62 patients (18.7\%) are 61 years or older
\end{itemize}

\subsection{Clinical Implications}

\begin{itemize}
\item \textbf{Healthcare Utilization:} Peak utilization in 41--60 age group suggests high prevalence of chronic conditions in middle age
\item \textbf{Preventive Care Opportunities:} Large young adult population (29.3\%) presents opportunities for preventive interventions
\item \textbf{Geriatric Considerations:} Nearly 1 in 5 patients requires age-specific geriatric care considerations
\item \textbf{Workforce Health:} High representation of working-age adults (31--60 years) indicates significant occupational health implications
\end{itemize}

\section{Age-Related Healthcare Patterns}

\subsection{Distribution Analysis by Quartiles}

\begin{table}[h!]
\centering
\caption{Population Quartile Analysis}
\begin{tabular}{l|c|c|c}
\toprule
\rowcolor{headercolor}
\textbf{Quartile} & \textbf{Cumulative \%} & \textbf{Age Range} & \textbf{Patient Count} \\
\midrule
First Quartile (Q1) & 25\% & 11--38 years & 83 patients \\
\rowcolor{altrowcolor}
Second Quartile (Q2) & 50\% & 39--49 years & 82 patients \\
Third Quartile (Q3) & 75\% & 50--58 years & 83 patients \\
\rowcolor{altrowcolor}
Fourth Quartile (Q4) & 100\% & 59--100 years & 83 patients \\
\bottomrule
\end{tabular}
\end{table}

\subsection{Age Group Risk Stratification}

\begin{table}[h!]
\centering
\caption{Risk Categories by Age Groups}
\begin{tabular}{l|c|l}
\toprule
\rowcolor{headercolor}
\textbf{Risk Category} & \textbf{Patient Count} & \textbf{Healthcare Considerations} \\
\midrule
Low Risk (21--40) & 97 & Preventive care, lifestyle counseling \\
\rowcolor{altrowcolor}
Moderate Risk (41--60) & 161 & Chronic disease screening, management \\
High Risk (61--80) & 58 & Comprehensive geriatric assessment \\
\rowcolor{altrowcolor}
Very High Risk (81+) & 4 & Intensive monitoring, palliative considerations \\
Special Population (0--20) & 11 & Pediatric/adolescent specialized care \\
\bottomrule
\end{tabular}
\end{table}

\section{Comparative Demographics}

\subsection{Age Distribution Characteristics}

\begin{itemize}
\item \textbf{Normal Distribution Tendency:} The age distribution shows characteristics of a right-skewed normal distribution
\item \textbf{Central Tendency:} Concentration around middle age (40--60 years) suggests typical healthcare-seeking behavior
\item \textbf{Minimal Extremes:} Very low representation at both pediatric (3.3\%) and very elderly (1.2\%) extremes
\item \textbf{Healthcare Demand Peak:} The 41--60 age group peak aligns with typical onset of chronic conditions
\end{itemize}

\subsection{Population Health Implications}

\begin{enumerate}
\item \textbf{Resource Allocation:} Healthcare resources should prioritize middle-aged patient care infrastructure
\item \textbf{Preventive Programs:} Strong focus needed on preventive care for the large 21--40 age group
\item \textbf{Chronic Disease Management:} Comprehensive programs required for the 48.7\% middle-aged population
\item \textbf{Geriatric Services:} Specialized geriatric services needed for 18.7\% elderly population
\item \textbf{Pediatric Considerations:} Limited but important pediatric care capacity requirements
\end{enumerate}

\section{Statistical Validation}

\begin{table}[h!]
\centering
\caption{Data Quality and Validation Metrics}
\begin{tabular}{l|c}
\toprule
\rowcolor{headercolor}
\textbf{Validation Metric} & \textbf{Result} \\
\midrule
Total Patient Count & 331 patients \\
\rowcolor{altrowcolor}
Data Completeness & 100\% (no missing age data) \\
Age Range Coverage & 11--100 years \\
\rowcolor{altrowcolor}
Largest Single Group & 41--50 years (25.1\%) \\
Smallest Non-zero Group & 91--100 years (0.3\%) \\
\rowcolor{altrowcolor}
Median Age Group & 41--50 years \\
\bottomrule
\end{tabular}
\end{table}

\section{Methodology and Data Source}

\subsection{Data Collection}
\begin{itemize}
\item \textbf{Source:} Medical dataset (521.xlsx)
\item \textbf{Sample Size:} 331 patients
\item \textbf{Age Grouping:} 10-year intervals for demographic analysis
\item \textbf{Data Validation:} Complete age information for all patients
\end{itemize}

\subsection{Statistical Methods}
\begin{itemize}
\item \textbf{Descriptive Statistics:} Frequency distribution and percentage calculations
\item \textbf{Categorical Analysis:} Age group consolidation for clinical relevance
\item \textbf{Cumulative Analysis:} Progressive percentage calculations for quartile identification
\item \textbf{Risk Stratification:} Age-based risk category assignments
\end{itemize}

\section{Conclusions}

The age distribution analysis reveals a patient population predominantly composed of middle-aged adults (41--60 years, 48.7\%), with significant representation of young adults (21--40 years, 29.3\%) and moderate elderly participation (61+ years, 18.7\%). This distribution pattern suggests a healthcare system primarily serving working-age adults with emerging chronic conditions, while maintaining capacity for preventive care in younger populations and specialized geriatric services for older patients.

The demographic profile indicates optimal resource allocation should focus on:
\begin{enumerate}
\item Chronic disease management for middle-aged patients
\item Preventive health programs for young adults
\item Specialized geriatric care for elderly patients
\item Limited but essential pediatric services
\end{enumerate}

\vspace{1cm}

\textbf{Analysis Date:} October 2025\\
\textbf{Dataset:} 521.xlsx medical records\\
\textbf{Font Standard:} Times New Roman (mathptmx)\\
\textbf{Formatting:} IEEE journal guidelines

\end{document}